\chapter{Nuclear Medicine \& Molecular Imaging}

\medterm{Annihilation Radiation}-- Pair of $511\text{ keV}$ gamma-ray photons produced when a positron emitted from a radioisotope encounters an electron and undergoes mutual annihilation, traveling in opposite directions ($180^\circ$ apart); forms the physical basis for Positron Emission Tomography (PET).

\medterm{Biological Half-Life}-- Time required for the human body to eliminate half the administered quantity of a drug or radiopharmaceutical through physiological excretory pathways (renal, fecal, pulmonary, hepatic), independent of radioactive decay.

\medterm{Collimator}-- Lead or tungsten shielding device mounted on a gamma camera or SPECT detector containing precision holes that permit only gamma rays traveling along specific geometric trajectories to reach the scintillation detector, preventing scattered photons from blurring the image.

\medterm{Cyclotron}-- Particle accelerator used in nuclear medicine radiopharmacies to bombard stable target nuclei with high-energy protons or deuterons to manufacture short-lived positron-emitting radionuclides such as Fluorine-18 ($^{18}\text{F}$), Carbon-11 ($^{11}\text{C}$), Nitrogen-13 ($^{13}\text{N}$), and Oxygen-15 ($^{15}\text{O}$).

\medterm{DaTscan (Ioflupane I-123)}-- Radiopharmaceutical ($^{123}\text{I-ioflupane}$) SPECT imaging agent that binds to presynaptic dopamine transporters (DaT) in the striatum; used to differentiate Essential Tremor (normal uptake) from Parkinsonian Syndromes (Parkinson disease, Multiple System Atrophy, Progressive Supranuclear Palsy showing reduced striatal uptake).

\medterm{DMSA Renal Scan}-- Nuclear medicine imaging procedure using Technetium-99m dimercaptosuccinic acid ($^{99m}\text{Tc-DMSA}$) which binds to cortical renal tubular cells; gold standard for evaluating static renal cortical mass, detecting cortical scarring from recurrent pyelonephritis, and assessing differential renal function.

\medterm{Dose Calibrator}-- Ionization chamber instrument utilized in radiopharmacies to measure the radioactivity (Bequerels or Curies) of radiopharmaceutical doses prior to administration to patients.

\medterm{DTPA Renal Scan}-- Dynamic renal scintigraphy using Technetium-99m diethylenetriaminepentaacetic acid ($^{99m}\text{Tc-DTPA}$), an agent excreted exclusively by glomerular filtration; used to measure Glomerular Filtration Rate (GFR) and assess renal arterial perfusion.

\medterm{Effective Half-Life}-- Net rate of clearance of a radiopharmaceutical from the body combining both physical radioactive decay half-life ($T_p$) and biological clearance half-life ($T_b$), calculated by $(T_p \times T_b) / (T_p + T_b)$.

\medterm{Fluorodeoxyglucose (18F-FDG)}-- Positron-emitting radiopharmaceutical and glucose analog ($^{18}\text{F-FDG}$) widely utilized in Positron Emission Tomography (PET) imaging. Pathophysiology: FDG is transported into cells via GLUT transporters and phosphorylated by hexokinase into FDG-6-phosphate; because FDG-6-phosphate cannot be further metabolized in the glycolytic pathway, it becomes metabolic trapped inside hypermetabolic cells. Clinical utility: Oncology (detection, staging, and therapy response monitoring of primary and metastatic malignancies, which exhibit high glucose utilization via the Warburg effect), Neurology (evaluation of Alzheimer's disease showing parietotemporal hypometabolism, and localization of epileptic foci), and Cardiology (assessment of myocardial viability).

\medterm{Fluorodeoxyglucose PET-CT (FDG-PET/CT)}-- Hybrid molecular imaging modality combining high-sensitivity Positron Emission Tomography (PET) tracking glucose metabolism with high-resolution Computed Tomography (CT) providing precise anatomical localization. Essential in modern clinical oncology for tumor staging (TNM), restaging following therapy, detecting occult recurrence, guiding biopsy of hypermetabolic tissue, and planning radiation therapy fields. Standardized Uptake Value (SUV) is a quantitative measurement of radiotracer uptake used to differentiate malignant from benign lesions.

\medterm{Gallium-68 PSMA (68Ga-PSMA)}-- Positron-emitting radiopharmaceutical targeting Prostate-Specific Membrane Antigen (PSMA), highly expressed on prostate adenocarcinoma cells; utilized in PET-CT for restaging prostate cancer with biochemical recurrence (elevated PSA) and detecting occult micrometastases.

\medterm{Gamma Camera}-- Core imaging detector in conventional nuclear medicine, designed to detect gamma rays emitted from radiopharmaceuticals administered to patients and convert them into two-dimensional projection images (scintigrams). Major components: lead collimator (discriminates gamma ray direction), sodium iodide (NaI) scintillation crystal (converts gamma photons into visible light flashes), photomultiplier tubes (PMTs—convert light into electrical signals), and positioning electronics. Used for planar imaging and rotated around the patient to perform Single-Photon Emission Computed Tomography (SPECT).

\medterm{Gastric Emptying Study}-- Diagnostic nuclear medicine procedure where a patient ingests a radiolabeled standardized meal (e.g., egg whites labeled with $^{99m}\text{Tc-sulfur colloid}$) and sequential gamma camera images are acquired to measure gastric retention percentage over 4 hours; gold standard for diagnosing gastroparesis.

\medterm{HIDA Scan (Cholescintigraphy)}-- Dynamic hepatobiliary nuclear imaging using Technetium-99m iminodiacetic acid analogs ($^{99m}\text{Tc-mebrofenin}$ or $^{99m}\text{Tc-disofenin}$); excreted by hepatocytes into bile, allowing visualization of gallbladder, biliary tree, and duodenum. Non-visualization of gallbladder at 4 hours indicates acute cholecystitis due to cystic duct obstruction.

\medterm{Indium-111 WBC Scan}-- Diagnostic imaging technique where patient leukocytes are harvested, labeled in vitro with Indium-111 oxine ($^{111}\text{In-WBC}$), and reinjected; highly specific for localizing occult intra-abdominal abscesses, osteomyelitis in complicated diabetic feet, and inflammatory bowel disease.

\medterm{Iodine-123 MIBG (123I-MIBG)}-- Radiopharmaceutical structural analog of norepinephrine concentrated in sympathetic adrenergic tissue; used in SPECT imaging to detect neuroblastoma, pheochromocytoma, and paraganglioma, and assess cardiac sympathetic innervation.

\medterm{Krypton-81m (81mKr)}-- Ultra-short half-life (13 seconds) inert radioactive gas generated on-site from a Rubidium-81 generator; inhaled to perform real-time continuous pulmonary ventilation scintigraphy paired with perfusion imaging.

\medterm{Lutetium-177 DOTATATE (177Lu-DOTATATE)}-- Targeted peptide receptor radionuclide therapy (PRRT) utilizing a beta-emitting radioisotope linked to a somatostatin analog; indicated for the treatment of somatostatin receptor-positive gastroenteropancreatic neuroendocrine tumors (GEP-NETs).

\medterm{Lymphoscintigraphy}-- Nuclear medicine procedure mapping lymphatic drainage pathways following intradermal injection of $^{99m}\text{Tc-sulfur colloid}$ or $^{99m}\text{Tc-tilmanocept}$ around a primary malignancy (melanoma, breast cancer) to identify and guide surgical biopsy of the Sentinel Lymph Node.

\medterm{MAG3 Renal Scan}-- Dynamic renal scintigraphy using Technetium-99m mercaptoacetyltriglycine ($^{99m}\text{Tc-MAG3}$), excreted predominantly by renal tubular secretion; ideal radiotracer for evaluating urinary tract obstruction (hydronephrosis) and renal dysfunction in patients with impaired renal function.

\medterm{Meckel Scintigraphy}-- Nuclear medicine scan following intravenous injection of Technetium-99m pertechnetate ($^{99m}\text{TcO}_4^-$), which selectively accumulates in ectopic gastric mucosa; gold standard diagnostic test for detecting Meckel's diverticulum presenting with lower gastrointestinal bleeding.

\medterm{Myocardial Perfusion Imaging (MPI)}-- Diagnostic nuclear cardiology test used to evaluate myocardial blood flow at rest and during physiological or pharmacological stress (adenosine, regadenoson, dobutamine). Radiopharmaceuticals used: Technetium-99m sestamibi ($^{99m}\text{Tc-sestamibi}$), Technetium-99m tetrofosmin, or Thallium-201 ($^{201}\text{Tl}$). Images captured using SPECT camera allow detection of inducible ischemia (reversible perfusion defect: normal at rest, reduced during stress—indicating significant coronary artery disease) versus myocardial infarction (fixed perfusion defect: reduced at both rest and stress—indicating scar tissue).

\medterm{Nuclear Medicine \& Molecular Imaging}-- Specialized branch of medical imaging and therapeutics that utilizes small amounts of radioactive substances (radiopharmaceuticals or radiotracers) to diagnose, stage, evaluate, and treat human diseases. Unlike anatomical imaging (plain X-ray, CT, MRI) which visualizes structure, nuclear medicine visualizes functional physiological processes at the cellular and molecular level (organ perfusion, cell metabolism, receptor density, bone turnover). Encompasses diagnostic modalities (SPECT, PET, planar scintigraphy) and targeted radionuclide therapy (Theranostics).

\medterm{Parathyroid Scintigraphy}-- Subtractive or dual-phase SPECT/CT imaging utilizing Technetium-99m sestamibi ($^{99m}\text{Tc-sestamibi}$) which washes out rapidly from normal thyroid tissue but is retained in hyperfunctioning parathyroid adenomas or hyperplasia, guiding minimally invasive parathyroidectomy.

\medterm{PET-MRI}-- Hybrid diagnostic imaging modality combining the high metabolic sensitivity of Positron Emission Tomography (PET) with the superior soft-tissue contrast and non-ionizing anatomical imaging of Magnetic Resonance Imaging (MRI).

\medterm{Positron Decay}-- Mode of radioactive decay ($\\beta^+$ decay) occurring in proton-rich unstable nuclei, where a proton converts into a neutron, emitting a positron ($e^+$) and a neutrino ($\nu$).

\medterm{Positron Emission Tomography (PET)}-- Advanced 3D tomographic molecular imaging technique based on the detection of annihilation photon pairs. When a positron-emitting radiotracer (e.g., Fluorine-18, Carbon-11, Gallium-68) decays, emitted positrons travel a short distance before colliding with an electron, resulting in an annihilation event that produces two $511\text{ keV}$ gamma ray photons traveling in opposite directions ($180^\circ$ apart). Detectors surrounding the patient register these simultaneous coincidence events to construct high-resolution functional images.

\medterm{Radioactive Iodine Therapy (131I)}-- Targeted radionuclide therapy utilizing Iodine-131 ($^{131}\text{I}$), a beta- and gamma-emitting isotope, to selectively destroy functioning thyroid tissue. Thyroid follicular cells actively trap circulating iodide via the sodium-iodide symporter (NIS). Emitted beta particles (mean range ~1--2 mm) deliver localized cytotoxic radiation while sparing adjacent neck structures. Clinical indications: definitive treatment of hyperthyroidism (Graves' disease, toxic multinodular goiter) and adjuvant ablation of residual normal thyroid tissue or metastatic disease following surgery for differentiated thyroid cancer (papillary and follicular carcinoma).

\medterm{Radioimmunoassay (RIA)}-- In vitro quantitative laboratory technique utilizing radiolabeled antigens ($^{125}\text{I}$) and specific antibodies to measure extremely low concentrations of hormones, drugs, and proteins in blood plasma.

\medterm{Radionuclide Bone Scan}-- Highly sensitive nuclear medicine imaging procedure used to evaluate bone turnover and osteoblastic activity throughout the skeleton. Radiopharmaceutical: Technetium-99m labeled diphosphonates ($^{99m}\text{Tc-MDP}$), which adsorb directly onto hydroxyapatite crystals in areas of active bone remodeling. Three-phase bone scan (flow phase, blood pool phase, delayed skeletal phase) differentiates osteomyelitis from cellulitis. Highly sensitive for detecting skeletal metastases (prostate, breast, lung cancer—appearing as multiple asymmetrical "hot spots"), stress fractures, occult trauma, and Paget disease of bone.

\medterm{Radiopharmaceutical}-- Pharmaceutical drug compound containing one or more radioactive isotopes (radionuclides) suitable for human administration, used for medical diagnosis or targeted therapy. Composed of two parts: a targeting molecule (antibody, peptide, metabolic substrate, or ligand) that binds to specific cellular receptors or physiological pathways, linked to a radionuclide that emits detectable radiation (gamma or positron for imaging; beta or alpha for therapy). Common radionuclides: Technetium-99m ($^{99m}\text{Tc}$), Fluorine-18 ($^{18}\text{F}$), Iodine-131 ($^{131}\text{I}$), Gallium-68 ($^{68}\text{Ga}$), Lutetium-177 ($^{177}\text{Lu}$), and Yttrium-90 ($^{90}\text{Y}$).

\medterm{Radium-223 Dichloride (223Ra)}-- Alpha-emitting targeted radiopharmaceutical that acts as a calcium mimetic, selectively targeting areas of high bone turnover; approved for improving overall survival in castration-resistant prostate cancer with symptomatic bone metastases.

\medterm{Samarium-153 Lexidronam (153Sm)}-- Beta-emitting radiopharmaceutical complexed with tetramethylene phosphonate, concentrating in osteoblastic bone lesions; used for palliative pain relief in widespread painful skeletal metastases.

\medterm{Schilling Test}-- Historical classic radioisotope absorption test utilizing Cobalt-57 labeled Vitamin B12 ($^{57}\text{Co-B12}$) to elucidate the underlying etiology of megaloblastic anemia (differentiating intrinsic factor deficiency / pernicious anemia from ileal malabsorption).

\medterm{Secular Equilibrium}-- Radioactive equilibrium established in a parent-daughter isotope pair when the physical half-life of the parent radionuclide is extremely long compared to that of the daughter (e.g., Radium-226 decaying to Radon-222).

\medterm{Single-Photon Emission Computed Tomography (SPECT)}-- 3D tomographic nuclear imaging technique where one or more gamma cameras rotate around the patient, acquiring multiple 2D projection images from different angles that are computer-reconstructed into cross-sectional axial, coronal, and sagittal slices. Utilizes single-photon gamma-emitting radiotracers (primarily Technetium-99m, Iodine-123, Indium-111). Frequently combined with CT (hybrid SPECT/CT) for precise anatomical registration. Used extensively in myocardial perfusion imaging, bone scans, brain SPECT (DaTscan for Parkinsonism), and neuroendocrine tumor localization.

\medterm{Standardized Uptake Value (SUV)}-- Semi-quantitative measurement of radiotracer concentration in a region of interest normalized to injected dose and patient body weight, calculated during PET imaging to differentiate malignant (high SUV) from benign lesions.

\medterm{Strontium-89 Chloride (89Sr)}-- Beta-emitting radioisotope and calcium analog that deposits in active osteoblastic bone lesions; utilized for long-term palliative relief of bone pain caused by osseous metastases.

\medterm{Technetium-99m (99mTc)}-- The most widely used diagnostic radionuclide in modern nuclear medicine, accounting for over 80\% of all nuclear imaging procedures globally. Ideal physical properties: emits pure $140\text{ keV}$ gamma rays (optimal energy for gamma camera detection with minimal radiation dose to patient), has a convenient physical half-life of 6.0 hours (short enough to minimize radiation exposure, long enough for complex imaging protocols), and decays by isomeric transition to Technetium-99 without emitting particulate (beta/alpha) radiation. Produced on-site in hospital radiopharmacies using a Molybdenum-99 / Technetium-99m generator ("moly cow").

\medterm{Technetium-99m HMPAO}-- Lipophilic radiopharmaceutical ($^{99m}\text{Tc-hexamethylpropyleneamine oxime}$) that crosses the blood-brain barrier and becomes intracellularly trapped in brain tissue in proportion to regional cerebral blood flow; utilized in SPECT brain imaging to confirm Brain Death.

\medterm{Technetium-99m MAA}-- Technetium-99m labeled Macroaggregated Albumin ($^{99m}\text{Tc-MAA}$), consisting of radiolabeled albumin particles ($20--90~\mu\text{m}$) that temporary lodge in pulmonary capillaries; used for pulmonary perfusion imaging and pre-transarterial radioembolization (TARE) planning.

\medterm{Technetium-99m Sestamibi}-- Lipophilic cationic radiopharmaceutical ($^{99m}\text{Tc-sestamibi}$) that passively diffuses into cells and concentrates in mitochondrial membranes; used extensively in myocardial perfusion imaging, parathyroid imaging, and breast imaging.

\medterm{Thallium-201 (201Tl)}-- Cyclotron-produced potassium analog radionuclide emitting low-energy X-rays; historically used for myocardial perfusion imaging and differentiating cerebral lymphoma from toxoplasmosis in immunocompromised patients.

\medterm{Theranostics}-- Emerging paradigm in nuclear medicine and oncology that combines targeted diagnostic imaging and targeted radionuclide therapy using the same molecular vector or paired radiotracers. A diagnostic radiotracer (e.g., Gallium-68 PSMA for prostate cancer, or $^{68}\text{Ga-DOTATATE}$ for neuroendocrine tumors) is first used with PET-CT to visualize, localize, and quantify target receptor expression on tumor cells. If positive, the same targeting vector is labeled with a therapeutic beta- or alpha-emitting radionuclide (e.g., Lutetium-177 PSMA, or $^{177}\text{Lu-DOTATATE}$) to deliver destructive ionizing radiation selectively to tumor sites while sparing healthy tissue.

\medterm{Thyroid Scintigraphy (123I / 99mTc)}-- Diagnostic nuclear imaging procedure evaluating the anatomical structure and functional activity of the thyroid gland. Radiotracers: Iodine-123 ($^{123}\text{I}$) or Technetium-99m pertechnetate ($^{99m}\text{TcO}_4^-$). Clinical indications: differential diagnosis of hyperthyroidism—diffuse homogeneous increased uptake (Graves' disease), toxic nodular goiter (multiple hot nodules), toxic adenoma (single hyperfunctioning "hot" nodule suppressing rest of gland), or subacute thyroiditis (suppressed, absent uptake); and evaluation of thyroid nodules ("cold" non-functioning nodules carry a ~15--20\% risk of malignancy and require fine-needle aspiration).

\medterm{Transient Equilibrium}-- Radioactive equilibrium condition occurring when the physical half-life of the parent radionuclide is slightly longer than that of the daughter radionuclide, as seen in Molybdenum-99 ($T_{1/2}=66\text{ h}$) decaying to Technetium-99m ($T_{1/2}=6\text{ h}$).

\medterm{Ventilation-Perfusion Scan (V/Q Scan)}-- Nuclear medicine study evaluating pulmonary ventilation and blood perfusion to diagnose pulmonary embolism (PE), particularly in patients with contraindications to CT Pulmonary Angiography (e.g., severe contrast allergy, renal failure, pregnancy). Ventilation phase: inhalation of Radioactive Xenon-133 gas or $^{99m}\text{Tc}$-Technegas aerosol. Perfusion phase: intravenous injection of Technetium-99m macroaggregated albumin ($^{99m}\text{Tc-MAA}$), which microembolizes in pulmonary capillaries. Classic PE finding: one or more segmental "mismatched" defects (normal ventilation with absent perfusion).

\medterm{Xenon-133 (133Xe)}-- Inhaled radioactive inert gas emitting $81\text{ keV}$ gamma rays; used in planar nuclear ventilation imaging to evaluate regional pulmonary airflow and trapped gas in chronic obstructive pulmonary disease.

\medterm{Yttrium-90 Microspheres (90Y)}-- Radioembolization selective internal radiation therapy (SIRT) utilizing glass or resin microspheres loaded with Yttrium-90 ($^{90}\text{Y}$), delivered via hepatic artery catheterization to treat primary hepatocellular carcinoma or unresectable liver metastases.

